%% Created by Maple 2022.1, Windows 10
%% Source Worksheet: lab_sols06
%% Generated: Thu May 09 11:30:56 BST 2024
\documentclass{article}
\usepackage{amssymb}
\usepackage{graphicx}
\usepackage{hyperref}
\usepackage{listings}
\usepackage{mathtools}
\usepackage{maple}
\usepackage[utf8]{inputenc}
\usepackage[svgnames]{xcolor}
\usepackage{amsmath}
\usepackage{breqn}
\usepackage{textcomp}
\begin{document}
\lstset{basicstyle=\ttfamily,breaklines=true,columns=flexible}
\pagestyle{empty}
\DefineParaStyle{Maple Bullet Item}
\DefineParaStyle{Maple Heading 1}
\DefineParaStyle{Maple Warning}
\DefineParaStyle{Maple Heading 4}
\DefineParaStyle{Maple Heading 2}
\DefineParaStyle{Maple Heading 3}
\DefineParaStyle{Maple Dash Item}
\DefineParaStyle{Maple Error}
\DefineParaStyle{Maple Title}
\DefineParaStyle{Maple Ordered List 1}
\DefineParaStyle{Maple Text Output}
\DefineParaStyle{Maple Ordered List 2}
\DefineParaStyle{Maple Ordered List 3}
\DefineParaStyle{Maple Normal}
\DefineParaStyle{Maple Ordered List 4}
\DefineParaStyle{Maple Ordered List 5}
\DefineCharStyle{Maple 2D Output}
\DefineCharStyle{Maple 2D Input}
\DefineCharStyle{Maple Maple Input}
\DefineCharStyle{Maple 2D Math}
\DefineCharStyle{Maple Hyperlink}
\begin{center}
\begin{Maple Title}
\underline{\textbf{More differentiation}}
\end{Maple Title}
\end{center}
\begin{Maple Normal}
\_\_\_\_\_\_\_\_\_\_\_\_\_\_\_\_\_\_\_\_\_\_\_\_\_\_\_\_\_\_\_\_\_\_\_\_\_\_\_\_\_\_\_\_\_\_\_\_\_\_\_\_\_\_\_\_\_\_\_\_\_\_\_\_\_\_\_\_\_\_\_\_\_\_\_\_\_\_\_\_\_
\end{Maple Normal}
\begin{Maple Heading 2}
\textbf{Exercise 1.1}
\end{Maple Heading 2}
\begin{Maple Normal}
Here are the graphs of 
$y =x^{n} {\mathrm e}^{-x}$ for 
$n =3,4$ and 
$5$:
\end{Maple Normal}
\begin{lstlisting}
> plot(x^3*exp(-x),x=0..8);
plot(x^4*exp(-x),x=0..8);
plot(x^5*exp(-x),x=0..8);
\end{lstlisting}
\begin{Maple Normal}
You can just about see that the peaks occur at 
$x =3$, 
$x =4$ and 
$x =5$, suggesting that for general 
$n$ 
\end{Maple Normal}
\begin{Maple Normal}
we should have a peak at 
$x =n$.  To check this, we must solve 
$\frac{\mathit{dy}}{\mathit{dx}}=0$.
\end{Maple Normal}
\begin{lstlisting}
> y := x^n * exp(-x);
\end{lstlisting}
\begin{lstlisting}
> simplify(diff(y,x));
\end{lstlisting}
\begin{lstlisting}
> solve(diff(y,x)=0,\{x\});
\end{lstlisting}
\begin{Maple Normal}
This shows that the peak does indeed occur at 
$x =n$.  To find the height of the peak, we must put
\end{Maple Normal}
\begin{Maple Normal}

$x =n$ in 
$y$:
\end{Maple Normal}
\begin{lstlisting}
> y[max] := subs(x=n,y);
\end{lstlisting}
\begin{Maple Normal}
It follows that 
$\frac{y}{y_{\max}}=(\frac{x}{n})^{n} {\mathrm e}^{n -x}$, so the maximum value of  
$(\frac{x}{n})^{n} {\mathrm e}^{n -x}$ is 
$1$.  We can plot these functions for different 
$n$ together as follows:
\end{Maple Normal}
\begin{lstlisting}
> plot([seq((x/n)^n*exp(n-x),n=1..8)],x=0..20);
\end{lstlisting}
\begin{Maple Normal}
\_\_\_\_\_\_\_\_\_\_\_\_\_\_\_\_\_\_\_\_\_\_\_\_\_\_\_\_\_\_\_\_\_\_\_\_\_\_\_\_\_\_\_\_\_\_\_\_\_\_\_\_\_\_\_\_\_\_\_\_\_\_\_\_\_\_\_\_\_\_\_\_\_\_\_\_\_\_\_\_\_
\end{Maple Normal}
\begin{Maple Heading 2}
\textbf{Exercise 1.2}
\end{Maple Heading 2}
\begin{lstlisting}
> p := (x) -> -10*x^6+156*x^5-945*x^4+2780*x^3-4080*x^2+2880*x;
\end{lstlisting}
\begin{lstlisting}
> plot(p(x),x=0..5);
\end{lstlisting}
\begin{Maple Normal}
We see from the picture that the maximum value is about 1000, occurring at about 
$x =3$. 
\end{Maple Normal}
\begin{Maple Normal}
There also seem to be two inflection points, where the curve flattens out but does not
\end{Maple Normal}
\begin{Maple Normal}
have a maximum or minimum.  To check this, we solve for 
$\mathit{p'} (x)=0$:
\end{Maple Normal}
\begin{lstlisting}
> p1 := diff(p(x),x);
\end{lstlisting}
\begin{lstlisting}
> solve(p1=0,x);
\end{lstlisting}
\begin{Maple Normal}
This shows that the critical points are at 
$x =1$, 
$x =3$ and 
$x =4$.  The numbers 1 and 4 are
\end{Maple Normal}
\begin{Maple Normal}
repeated because they are double roots of 
$\mathit{p'} (x)$, and therefore inflection points of 
$p (x)$.
\end{Maple Normal}
\begin{Maple Normal}
To check this, we differentiate again:
\end{Maple Normal}
\begin{lstlisting}
> p2 := diff(p(x),x,x);
\end{lstlisting}
\begin{lstlisting}
> subs(x=1,p2);
\end{lstlisting}
\begin{lstlisting}
> subs(x=3,p2);
\end{lstlisting}
\begin{lstlisting}
> subs(x=4,p2);
\end{lstlisting}
\begin{Maple Normal}
As 
$\mathit{p'} (3)=0$ and 
$\mathit{p''} (3)<0$ we see that 
$x =3$ is a local maximum; by looking at the graph we
\end{Maple Normal}
\begin{Maple Normal}
see that it is a global maximum.  The maximum value of 
$p (x)$ is thus given by 
$p (3)$:
\end{Maple Normal}
\begin{lstlisting}
> p(3);
\end{lstlisting}
\begin{Maple Normal}
The values of 
$p (x)$ at the inflection points are strictly smaller, as expected:
\end{Maple Normal}
\begin{lstlisting}
> p(1),p(4);
\end{lstlisting}
\begin{Maple Normal}
\_\_\_\_\_\_\_\_\_\_\_\_\_\_\_\_\_\_\_\_\_\_\_\_\_\_\_\_\_\_\_\_\_\_\_\_\_\_\_\_\_\_\_\_\_\_\_\_\_\_\_\_\_\_\_\_\_\_\_\_\_\_\_\_\_\_\_\_\_\_\_\_\_\_\_\_\_\_\_\_\_
\end{Maple Normal}
\begin{Maple Heading 2}
\textbf{Exercise 2.1}
\end{Maple Heading 2}
\begin{lstlisting}
> restart;
\end{lstlisting}
\begin{lstlisting}
> u := x*sin(x^2+y^2)-y*cos(x^2+y^2);
\end{lstlisting}
\begin{lstlisting}
> with(plots):
\end{lstlisting}
\begin{lstlisting}
> implicitplot(u=0,x=-5..5,y=-5..5,grid=[100,100]);
\end{lstlisting}
\begin{lstlisting}
> slope1 := implicitdiff(u=0,y,x);
\end{lstlisting}
\begin{Maple Normal}
The answer can be rewritten more nicely using the original relation
\end{Maple Normal}
\begin{center}
 
$x \sin (x^{2}+y^{2})=y \cos (x^{2}+y^{2})$, \end{center}
\begin{Maple Normal}

\end{Maple Normal}
\begin{Maple Normal}
which allows us to rewrite 
$\sin (x^{2}+y^{2})$ as 
$\frac{y \cos (x^{2}+y^{2})}{x}$, after which we can cancel 
\end{Maple Normal}
\begin{Maple Normal}
the cos terms.
\end{Maple Normal}
\begin{lstlisting}
> subs(sin(x^2+y^2) = (y/x)*cos(x^2+y^2),slope1);
\end{lstlisting}
\begin{lstlisting}
> simplify(%);
\end{lstlisting}
\begin{lstlisting}
> slope1 := %;
\end{lstlisting}
\begin{Maple Normal}

\end{Maple Normal}
\begin{Maple Normal}
We now write the curve parametrically in terms of 
$t$:
\end{Maple Normal}
\begin{lstlisting}
> xt := t * cos(t^2); yt := t * sin(t^2);
\end{lstlisting}
\begin{Maple Normal}
To check that this really is the same curve as before, we put 
$x =t \cos (t^{2})$ and 
$y =t \sin (t^{2})$ in 
$u$ and make sure that we get zero:
\end{Maple Normal}
\begin{lstlisting}
> subs(x=xt,y=yt,u);
\end{lstlisting}
\begin{lstlisting}
> simplify(subs(x=xt,y=yt,u));
\end{lstlisting}
\begin{Maple Normal}
We can also check graphically:
\end{Maple Normal}
\begin{lstlisting}
> plot([xt,yt,t=-5..5]);
\end{lstlisting}
\begin{Maple Normal}
We now calculate 
$\frac{\mathit{dy}}{\mathit{dx}}$ again from the parametric representation:
\end{Maple Normal}
\begin{lstlisting}
> diff(xt,t);
\end{lstlisting}
\begin{lstlisting}
> diff(yt,t);
\end{lstlisting}
\begin{lstlisting}
> slope2 := diff(yt,t)/diff(xt,t);
\end{lstlisting}
\begin{Maple Normal}
This should be the same as what we get by putting 
$x =t \cos (t^{2})$ and 
$y =t \sin (t^{2})$ in our earlier answer:
\end{Maple Normal}
\begin{lstlisting}
> simplify(subs(x=xt,y=yt,slope1));
\end{lstlisting}
\begin{Maple Normal}
\_\_\_\_\_\_\_\_\_\_\_\_\_\_\_\_\_\_\_\_\_\_\_\_\_\_\_\_\_\_\_\_\_\_\_\_\_\_\_\_\_\_\_\_\_\_\_\_\_\_\_\_\_\_\_\_\_\_\_\_\_\_\_\_\_\_\_\_\_\_\_\_\_\_\_\_\_\_\_\_\_
\end{Maple Normal}
\begin{Maple Heading 2}
\textbf{Exercise 2.2}
\end{Maple Heading 2}
\begin{lstlisting}
> restart;
\end{lstlisting}
\begin{Maple Normal}
We start by entering the definitions given in the question:
\end{Maple Normal}
\begin{lstlisting}
> u := (x^2+y^2)^2 + 85*(x^2+y^2) - 500 + 18*x*(3*y^2-x^2);
\end{lstlisting}
\begin{lstlisting}
> xt := 6*cos(t) + 8*cos(t)^2 - 4;
yt := 2*sin(t)*(3-4*cos(t));
\end{lstlisting}
\begin{Maple Normal}
We now plot the curve 
$u =0$ and the curve given parametrically in terms of 
$t$:
\end{Maple Normal}
\begin{lstlisting}
> with(plots):
\end{lstlisting}
\begin{lstlisting}
> implicitplot(u=0,x=-10..10,y=-10..10,grid=[100,100]);
\end{lstlisting}
\begin{lstlisting}
> plot([xt,yt,t=0..2*Pi]);
\end{lstlisting}
\begin{Maple Normal}
The curves certainly look the same.  This can be tested as follows:
\end{Maple Normal}
\begin{lstlisting}
> simplify(subs(x=xt,y=yt,u));
\end{lstlisting}
\begin{Maple Normal}
We next find 
$\frac{\mathit{dy}}{\mathit{dx}}$ by implicit differentiation:
\end{Maple Normal}
\begin{lstlisting}
> slope1 := simplify(implicitdiff(u=0,y,x));
\end{lstlisting}
\begin{Maple Normal}
Alternatively, we can find 
$\frac{\mathit{dy}}{\mathit{dx}}$ from the parametric representation:
\end{Maple Normal}
\begin{lstlisting}
> slope2 := simplify(diff(yt,t)/diff(xt,t));
\end{lstlisting}
\begin{Maple Normal}
This is the same as what we get by rewriting 
$\mathit{slope1}$ in terms of 
$t$:
\end{Maple Normal}
\begin{lstlisting}
> simplify(subs(x=xt,y=yt,slope1));
\end{lstlisting}
\begin{Maple Normal}
\_\_\_\_\_\_\_\_\_\_\_\_\_\_\_\_\_\_\_\_\_\_\_\_\_\_\_\_\_\_\_\_\_\_\_\_\_\_\_\_\_\_\_\_\_\_\_\_\_\_\_\_\_\_\_\_\_\_\_\_\_\_\_\_\_\_\_\_\_\_\_\_\_\_\_\_\_\_\_\_\_
\end{Maple Normal}
\begin{Maple Heading 2}
\textbf{Exercise 3.1}
\end{Maple Heading 2}
\begin{lstlisting}
> r := (n) -> diff(x^n*ln(x)/n!,x$n);
\end{lstlisting}
\begin{lstlisting}
> seq(r(n),n=1..10);
\end{lstlisting}
\begin{lstlisting}
> seq(r(n)-r(n-1),n=2..10);
\end{lstlisting}
\begin{Maple Normal}
We see from this that 
$r (n)-r (n -1)=\frac{1}{n}$.  We start with 
$r (1)=\ln (x)+1$, and then
\end{Maple Normal}
\begin{Maple Normal}
     
$r (2)=r (1)+r (2)-r (1)$
$=r (1)+\frac{1}{2}$
$=\ln (x)+1+\frac{1}{2}$
\end{Maple Normal}
\begin{Maple Normal}
     
$r (3)=r (2)+r (3)-r (2)$
$=r (2)+\frac{1}{3}$
$=\ln (x)+1+\frac{1}{2}+\frac{1}{3}$
\end{Maple Normal}
\begin{Maple Normal}
     
$r (4)=r (3)+r (4)-r (3)$
$=r (3)+\frac{1}{4}$
$=\ln (x)+1+\frac{1}{2}+\frac{1}{3}+\frac{1}{4}$
\end{Maple Normal}
\begin{Maple Normal}

\end{Maple Normal}
\begin{Maple Normal}
and so on.  In general, we have
\end{Maple Normal}
\begin{Maple Normal}

\end{Maple Normal}
\begin{center}

$r (n)=\ln (x)+1+\frac{1}{2}+\mathit{...} +\frac{1}{n}$
$=\ln (x)+\moverset{n}{\munderset{k =1}{\sum}}\frac{1}{k}$\end{center}
\begin{Maple Normal}

\end{Maple Normal}
\begin{Maple Normal}
Maple knows about a constant called 
$\mathrm{gamma}$ (Euler's constant, approximately 0.577) and a special 
\end{Maple Normal}
\begin{Maple Normal}
function called 
$\Psi$ with the property that 
$\moverset{n}{\munderset{k =1}{\sum}}\frac{1}{k}=\Psi (n +1)+\mathrm{gamma}$ for any positive integer 
$n$.
\end{Maple Normal}
\begin{Maple Normal}
(You can enter \textcolor[RGB]{255,0,0}{\textbf{?Psi}} to find out more.)  Using this, we could also write
\end{Maple Normal}
\begin{Maple Normal}

\end{Maple Normal}
\begin{center}
$\displaystyle r (n)=\ln (x)+\Psi (n +1)+\mathrm{gamma}$\end{center}
\begin{Maple Normal}
We can check this as follows:
\end{Maple Normal}
\begin{Maple Normal}

\end{Maple Normal}
\begin{lstlisting}
> seq(r(n),n=1..8);
\end{lstlisting}
\begin{lstlisting}
> seq(ln(x)+Psi(n+1)+gamma,n=1..8);
\end{lstlisting}
\begin{lstlisting}
> unassign('r');
\end{lstlisting}
\begin{Maple Normal}

\end{Maple Normal}
\begin{Maple Normal}
We can prove that the formula is valid by induction, but we will not give the details here. 
\end{Maple Normal}
\begin{Maple Normal}
\_\_\_\_\_\_\_\_\_\_\_\_\_\_\_\_\_\_\_\_\_\_\_\_\_\_\_\_\_\_\_\_\_\_\_\_\_\_\_\_\_\_\_\_\_\_\_\_\_\_\_\_\_\_\_\_\_\_\_\_\_\_\_\_\_\_\_\_\_\_\_\_\_\_\_\_\_\_\_\_\_
\end{Maple Normal}
\begin{Maple Heading 2}
\textbf{Exercise 3.2}
\end{Maple Heading 2}
\begin{lstlisting}
> y := t^2*exp(t);
\end{lstlisting}
\begin{lstlisting}
> z := simplify(
     diff(y,t,t,t) + 
 a * diff(y,t,t) + 
 b * diff(y,t) +
 c *y
); 
\end{lstlisting}
\begin{Maple Normal}
We want this to be zero for all 
$t$, so the coefficients of the individual powers of 
$t$ must all be
\end{Maple Normal}
\begin{Maple Normal}
zero.  To find these coefficients, we use the \textcolor[RGB]{255,0,0}{\textbf{collect}} function:
\end{Maple Normal}
\begin{lstlisting}
> collect(z,t);
\end{lstlisting}
\begin{Maple Normal}
We must therefore find 
$a ,b$ and 
$c$ such that 
\end{Maple Normal}
\begin{center}

$1+a +b +c =6+4 a +2 b$
$=6+2 a$
$=0$\end{center}
\begin{Maple Normal}

\end{Maple Normal}
\begin{lstlisting}
> solve(\{1+a+b+c=0,6+4*a+2*b=0,6+2*a=0\},\{a,b,c\});
\end{lstlisting}
\begin{Maple Normal}
The conclusion is that 
$\mathit{y'''} -3 \mathit{y''} +3 \mathit{y'} -y = 0.0$
\end{Maple Normal}
\begin{Maple Normal}

\end{Maple Normal}
\begin{Maple Normal}
You might think that we could do this more quickly by just solving 
$z =0$.  This does not work,
\end{Maple Normal}
\begin{Maple Normal}
because it does not capture the fact that 
$a ,b$ and 
$c$ are supposed to be constants.  Maple
\end{Maple Normal}
\begin{Maple Normal}
gives us the following answer, in which 
$b$ and 
$c$ can be anything, but 
$a$ depends on 
$t$:
\end{Maple Normal}
\begin{lstlisting}
> solve(z=0,\{a,b,c\});
\end{lstlisting}
\begin{Maple Normal}
A correct approach along these lines is to replace the equation \textcolor[RGB]{255,0,0}{\textbf{z=0}} by the expression \textcolor[RGB]{255,0,0}{\textbf{identity(z=0,t)}} to indicate that the equation is supposed to hold for all 
$t$. 
\end{Maple Normal}
\begin{Maple Normal}
This is probably the best approach, if you can remember the syntax.
\end{Maple Normal}
\begin{lstlisting}
> solve(identity(z=0,t),\{a,b,c\});
\end{lstlisting}
\begin{lstlisting}
> 
\end{lstlisting}
\begin{Maple Normal}
\_\_\_\_\_\_\_\_\_\_\_\_\_\_\_\_\_\_\_\_\_\_\_\_\_\_\_\_\_\_\_\_\_\_\_\_\_\_\_\_\_\_\_\_\_\_\_\_\_\_\_\_\_\_\_\_\_\_\_\_\_\_\_\_\_\_\_\_\_\_\_\_\_\_\_\_\_\_\_\_\_
\end{Maple Normal}
\begin{Maple Heading 2}
\textbf{Exercise 3.3}
\end{Maple Heading 2}
\begin{lstlisting}
> p := (n) -> sort(expand(exp(x^2)*diff(exp(-x^2),x$n)));
\end{lstlisting}
\begin{lstlisting}
> p(2);
\end{lstlisting}
\begin{lstlisting}
> p(3);
\end{lstlisting}
\begin{lstlisting}
> p(4);
\end{lstlisting}
\begin{lstlisting}
> seq(print(p(n)),n=1..10);
\end{lstlisting}
\begin{Maple Normal}
We first look at the leading term of  
$p (n)$.  The leading term of 
$p (1)$ is 
$-2 x$, the leading term of 
$p (2)$ is 
$4 x^{2}$, the leading term of 
$p (3)$ is 
$-8 x^{3}$ and so on.  The signs alternate, the constant is 
$2^{n}$, and the power of 
$x$ is 
$x^{n}$.  More succinctly, the leading term in 
$p (n)$ is 
$(-2 x)^{n}$.
\end{Maple Normal}
\begin{Maple Normal}

\end{Maple Normal}
\begin{Maple Normal}
Next, note that when 
$n$ is even we only get even powers of 
$x$, and when 
$n$ is odd we only get odd powers of 
$x$.  For example, 
$p (7)$ involves only 
$x ,x^{3},x^{5}$ and 
$x^{7}$, whereas 
$p (6)$ involves 
$x^{2},x^{4}$ and 
$x^{6}$ (and a constant term, which we can think of as a multiple of 
$x^{0}$).
\end{Maple Normal}
\begin{Maple Normal}

\end{Maple Normal}
\begin{Maple Normal}
Now look at the last term in 
$p (n)$.  When 
$n$ is even, the last term is a constant, but when 
$n$ is odd, it is a multiple of 
$x$.  It is best to consider these separately, starting with the even case, where we may write 
$n =2 m$.
\end{Maple Normal}
\begin{lstlisting}
> seq(print(p(2*m)),m=1..6);
\end{lstlisting}
\begin{Maple Normal}
We see that the last term is 
$(-1)^{m}$ times a constant, with the sequence of constants being 
$2,12,120,1680,302040,665280$ and so on.  If we enter this in the \href{http://www.research.att.com/~njas/sequences/}{Online Encyclopedia of Integer Sequences}  we get an answer including the line 
\end{Maple Normal}
\begin{Maple Normal}

\end{Maple Normal}
\begin{center}
\textbf{Name:  }    Quadruple factorial numbers: (2n)!/n!.\end{center}
\begin{Maple Normal}

\end{Maple Normal}
\begin{Maple Normal}
We need to be a little careful with this formula.  The encyclopedia assumes that the numbers we entered
\end{Maple Normal}
\begin{Maple Normal}
correspond to 
$n =1,2,3,\mathit{...}$ and gives the formula 
$\frac{(2 n)!}{n !}$.  In fact our numbers correspond to 
$m =1,2,3,\mathit{...}$ (where 
$n =2 m$) so the right formula for the constant term in 
$p (2 m)$ is 
$\frac{(-1)^{m} (2 m)!}{m !}$.
\end{Maple Normal}
\begin{Maple Normal}
Of course this kind of experimental approach does not really prove that the formula is correct, but it is very suggestive.
\end{Maple Normal}
\begin{Maple Normal}

\end{Maple Normal}
\begin{Maple Normal}
We now look at the case where 
$n$ is odd, say 
$n =2 m -1$:
\end{Maple Normal}
\begin{lstlisting}
> seq(print(p(2*m-1)),m=1..5);
\end{lstlisting}
\begin{Maple Normal}
The final terms are the same as before, but multiplied by 
$x$.  In other words, the last term in 
$p (2 m -1)$ 
\end{Maple Normal}
\begin{Maple Normal}
is 
$\frac{(-1)^{m} (2 m)! x}{m !}$. 
\end{Maple Normal}
\begin{Maple Normal}

\end{Maple Normal}
\begin{Maple Normal}
We can now predict that 
$p (12)$ should start with 
$(-2 x)^{12}=4096 x^{12}$, and that it should contain multiples of 
$x^{10},x^{8},x^{6},x^{4},x^{2}$ and a constant term.  This is the case where 
$n =2 m$ and 
$m =6$, so the constant term is 
$\frac{(-1)^{6}\cdot 12!}{6!}=665280$.  We check this as follows:
\end{Maple Normal}
\begin{lstlisting}
> p(12);
\end{lstlisting}
\begin{Maple Normal}
We now compare 
$p (n)$ with the Hermite polynomial 
$H_{n}(x)$, entered in Maple as \textcolor[RGB]{255,0,0}{\textbf{simplify(HermiteH(n,x))}}:
\end{Maple Normal}
\begin{lstlisting}
> q := (n) -> simplify(HermiteH(n,x));
\end{lstlisting}
\begin{lstlisting}
> seq(print([p(n),q(n)]),n=1..6);
\end{lstlisting}
\begin{Maple Normal}
We see that 
$q (n)$ is the same as 
$p (n)$ when 
$n$ is even, and the same as 
$-p (n)$ when 
$n$ is odd.  In other words, we have 
$q (n)=(-1)^{n} p (n)$.
\end{Maple Normal}
\begin{lstlisting}
> 
\end{lstlisting}
\begin{lstlisting}
> 
\end{lstlisting}
\end{document}
